\documentclass[10pt,a4paper]{amsart}
\usepackage[margin=19mm]{geometry}
\usepackage{amsmath,amssymb,mathtools}
\usepackage[scr=rsfs,cal=euler]{mathalfa}
\usepackage{xfrac}
\usepackage[unicode,hidelinks,bookmarks=false]{hyperref}
\newcommand{\ie}{i.e.\ }
\newcommand{\cf}{cf.\ }
\newcommand{\resp}{respectively\ }
\newcommand{\Subsection}{\subsection}
\providecommand{\nobreakdash}{\nobreak\hbox{-}}
\setlength{\emergencystretch}{2em}
\multlinegap0pt
\usepackage{bm}
\usepackage{bbm}
\usepackage{dsfont}
\usepackage{xspace}
\usepackage{cancel}
\usepackage{mathtools}
\usepackage{amsthm}
\makeatletter
\@ifundefined{theorem}{\newtheorem{theorem}{Theorem}}{}
\@ifundefined{lemma}{\newtheorem{lemma}[theorem]{Lemma}}{}
\@ifundefined{proposition}{\newtheorem{proposition}[theorem]{Proposition}}{}
\makeatother
\DeclareMathOperator{\tr}{tr}
\def\O{{\mathcal O}}
\def\S{{\mathcal S}}
\def\U{{\mathcal U}}
\def\V{{\mathcal V}}
\def\ol{\overline}
\def\vep{\varepsilon}
\newcommand{\Jop}{\overline{J}^{2,+}}
\newcommand{\R}{\mathbb{R}}
\title[Superposition Property in Disjoint Variables for the Infinit]{Superposition Property in Disjoint Variables for the Infinity Laplace Equation}
\author{Qing Liu, Juan J. Manfredi,, Xiaodan Zhou}
\date{Source: arXiv:2508.17106v2 (CC BY 4.0)}
\begin{document}
\maketitle

\noindent\emph{Source and licence.} Superposition Property in Disjoint Variables for the Infinity Laplace Equation. Version arXiv:2508.17106v2 (CC BY 4.0), distributed under the Creative Commons Attribution 4.0 International licence, as recorded in the arXiv licence metadata for this exact version and shown on the abstract page https://arxiv.org/abs/2508.17106v2. Full rights notice: licenses/arxiv-2508.17106-CC-BY-4.0.txt.\par\smallskip

\noindent\textit{Editorial scope.} Counted unit: the Proposition whose source label is \texttt{Lapsup} in the article (source0.tex lines 417--422, source label \texttt{Lapsup}), delivered together with the article's own complete original proof of it (source0.tex lines 425--505, one \texttt{proof} environment of 81 lines). No mathematics is written, repaired or re-derived: statement and proof are the article's own text line for line. Required context retained uncounted: TOS (lines 266--282). Every definition, display and label of those lines is retained unchanged. Omitted from this package and not redistributed: 1--265, 283--416, 423--424, 506--853. Nothing inside a retained excerpt is omitted or altered: every retained body is delivered exactly as the article's lines are written, so no omission receipt of any kind accompanies this package. Citations: the retained excerpts carry no citation command at all, so no bibliography entry of the article is transcribed and no reference is redistributed. Reference adaptation: none was needed and none was made; every reference inside the delivered unit resolves inside it, so no unresolved reference or citation is produced. Printed slips kept literal: no printed word, symbol, hypothesis or number of the counted statement or of its proof was altered. Layout and class: the article's own class is \texttt{amsart} with packages including inputenc, amsfonts, amsmath, amssymb, amsthm, enumerate, comment, esint, xcolor, geometry, soul; the wrapper is the lane's pinned \texttt{amsart} editorial wrapper at 10pt on A4, which restates the theorem-like environments the retained text needs and adds the article's own macro definitions that the retained text needs (\texttt{\textbackslash Jop}, \texttt{\textbackslash O}, \texttt{\textbackslash R}, \texttt{\textbackslash S}, \texttt{\textbackslash U}, \texttt{\textbackslash V}, \texttt{\textbackslash ol}, \texttt{\textbackslash tr}, \texttt{\textbackslash vep}), with extra wrapper packages (bm, bbm, dsfont, xspace, cancel, mathtools). The delivered numbering is the unit's own. Assets: no retained span contains a figure environment, a graphics-inclusion command, a PDF-page inclusion, a table or a TikZ picture; no image file is copied into this package, the package declares no asset dependency and no image file is redistributed with it. Licence: version arXiv:2508.17106v2 of this article is distributed under the Creative Commons Attribution 4.0 International licence, as recorded in the arXiv licence metadata for this exact version and shown on the abstract page https://arxiv.org/abs/2508.17106v2. A full rights notice is delivered at licenses/arxiv-2508.17106-CC-BY-4.0.txt. Characters: every retained excerpt is pure ASCII, so the delivered engine, XeLaTeX with its default Latin Modern text font, reports no missing character.
\begin{lemma}[Theorem on sums]\label{TOS}
Let $\U\subset \R^n$ and $\V\subset \R^m$ be two domains. Let $u:\U\to \R$ and $v:\V \to \R$ be upper semicontinuous, and $\varphi\in C^2(\U\times \V)$. Assume that \[ (x, y)\mapsto u(x)+v(y)-\varphi(x,y) \] reaches a maximum at $(x_0,y_0)\in \U\times \V$. For any $\varepsilon>0$, there exist $(\xi, X)\in \overline{J}^{2,+}u(x_0)$ and $(\eta, Y)\in \overline{J}^{2,+}v(y_0)$ such that
\[
\xi=\nabla_x \varphi(x_0,y_0),\quad\eta=\nabla_y \varphi(x_0,y_0),    
\]
and
\[
-\left(\frac{1}{\varepsilon}+\|Z\|\right)\begin{pmatrix}
I_n & 0 \\
0 & I_m 
\end{pmatrix}\le \begin{pmatrix}
X & 0 \\
0 & Y
\end{pmatrix}\le Z+\varepsilon Z^2,
\]
where $Z=\nabla^2 \varphi(x_0, y_0)\in \S^{n+m}$. 
\end{lemma}

\begin{proposition}[Superposition for Laplacian]\label{Lapsup}
    Let $\Omega\subset \mathbb{R}^n$ be a domain. Assume that $u, v$  are viscosity subsolutions (resp., supersolutions, solutions) to 
    \[
    -\Delta u=0 \quad \text{in $\Omega$}.
    \] Then $u+v$ is a viscosity subsolution (resp., supersolution, solution) to the same equation. 
\end{proposition}

\begin{proof}
Let $\varphi\in C^2(\Omega)$ and $\O\subset\subset \Omega$ be an open and bounded set. 
Assume that $u(x)+v(x)-\varphi(x)$ reaches a maximum in $\overline{\O}$ at $x=x_0$. Without changing the value of $-\Delta \varphi(x_0)$, we can add $|x-x_0|^4$ to $\varphi$ and assume that $x_0$ is the unique maximizer of $u+v-\varphi$ in $\ol{\O}$.   Suppose that
\[
\max_{x\in \overline{\O}}(u(x)+v(x)-\varphi(x))=u(x_0)+v(x_0)-\varphi(x_0)=\delta.
\]
Let $\varepsilon>0$. Consider $\Phi_\varepsilon:\overline{\O}\times\overline{\O}\to \R$ defined as
\[
\Phi_\varepsilon(x,y)=u(x)-\varphi(x)+v(y)-\frac{|x-y|^2}{2\varepsilon}.
\]
Then there exists $(x_\varepsilon, y_\varepsilon)\in \overline{\O}\times \overline{\O}$ such that
\[
\max_{\overline{\O}\times\overline{\O}}\Phi_\varepsilon(x,y)=\Phi(x_\varepsilon, y_\varepsilon)\ge \delta,
\]
which implies
\[
\frac{|x_\varepsilon-y_\varepsilon|^2}{2\varepsilon}\le u(x_\varepsilon)-\varphi(x_\varepsilon)+v(y_\varepsilon)-\delta.
\]
Due to the boundedness of $u$ and $v$ above, we therefore have $|x_\vep-y_\vep|\to 0$ as $\vep\to 0$. Since $\overline{\O}$ is compact, there exist subsequences such that $x_\varepsilon\to \hat{x}$ and $y_\varepsilon\to \hat{x}$ for some $\hat{x}\in \overline{\O}$. We still index the subsequences by $\vep$ for simplicity. 

By the upper semicontinuity of $u$ and $v$, we let $\vep\to 0$ to get
\[
\limsup_{\varepsilon\to 0}\frac{|x_\varepsilon-y_\varepsilon|^2}{2\varepsilon}\le u(\hat{x})-\varphi(\hat{x})+v(\hat{x})-\delta,
\]
which implies $u(\hat{x})-\varphi(\hat{x})+v(\hat{x})\geq \delta$. 
Since $u+v-\varphi$ achieves a unique maximum in $\ol{\O}$ at $x_0$, we therefore have $x_0=\hat{x}$. 

%and hence $\hat{x}=\hat{y}$. We deduce that $M\le 0$ and 
%\[
%\lim_{\varepsilon\to 0}\frac{|x_\varepsilon-%y_\varepsilon|^2}{2\varepsilon}=0,
%\]
%and 
%\[
%u(\hat{x})-\varphi(\hat{x})+v(\hat{x})=\delta=\max_{\overline{\O}} (u(x)+v(x)-\varphi(x)).
%\]
%\textcolor{blue}{Here $\hat{x}$ can be in $\partial \Omega$.}

Applying Lemma \ref{TOS} to 
\[
(x, y)\mapsto (u(x)-\varphi(x))+v(y)-\frac{|x-y|^2}{2\varepsilon},
\] 
we obtain %$X_\varepsilon,Y_\varepsilon\in S^n$ such that
\[
(\xi_\varepsilon, X_\varepsilon)\in \Jop(u-\varphi)(x_\varepsilon)\quad\text{ and } (\eta_\varepsilon, Y_\varepsilon)\in \Jop(v)(y_\varepsilon)\]
with
\[
\xi_\varepsilon=\frac{x_\varepsilon-y_\varepsilon}{\varepsilon}, \quad \eta_\varepsilon=\frac{y_\varepsilon-x_\varepsilon}{\varepsilon}
\]
and
\[
\begin{pmatrix}
X_\varepsilon & 0 \\
0 & Y_\varepsilon 
\end{pmatrix}\le \frac{3}{\varepsilon}\begin{pmatrix}
I & -I \\
-I & I 
\end{pmatrix}.
\]
Let $\zeta\in \mathbb{R}^n$ be an arbitrary vector. Multiplying the above inequality by $(\zeta, \zeta)$ from the left and its transpose $(\zeta, \zeta)^T$ from the right gives 
\begin{equation}\label{negH}
\langle(X_\varepsilon+Y_\varepsilon)\zeta, \zeta\rangle \le 0. 
\end{equation}
Hence, $\tr(X_\varepsilon)+\tr(Y_\varepsilon)\le 0.$
On the other hand, since $\varphi\in C^2(\Omega)$, we can easily verify that 
\[
(\xi_\varepsilon+\nabla \varphi(x_\varepsilon), \ X_\varepsilon+\nabla^2\varphi(x_\varepsilon))\in \Jop u(x_\varepsilon).
\]
By the fact that $u,v$ are viscosity subsolutions to $-\Delta u=0$, we have
\[
-\tr(X_\varepsilon+\nabla^2\varphi(x_\varepsilon))\le 0
\]
and 
\[
-\tr(Y_\varepsilon)\le 0.
\]
Adding the above inequalities and applying $\tr(X_\varepsilon)+\tr(Y_\varepsilon)\le 0$, we get
\[
-\Delta \varphi(x_\vep)=-\tr(\nabla^2\varphi(x_\varepsilon))\le 0
\]
 We conclude $-\Delta \varphi(x_0)\leq 0$ by letting $\varepsilon\to 0$. The proof is thus complete. 
\end{proof}

\end{document}
